%% The following is a directive for TeXShop to indicate the main file
%%!TEX root = diss.tex


%========================================
\chapter{Groups for walks with small steps}
%========================================

Consider $g_i(x_i) = X_i$ that fixes $S$ where $S$ is the Laurent step polynomial. We either have that $g_i$ is trivial or
\[
	g_i(x_i) = \ol{x}_i\frac{S_i^{-1}}{S_i^1},\quad g_i^2(x_i) = x_i
\]
whenever $S_i^{\pm 1} \neq 0$. These bi-rational transformations generate a group $G_S$
\[
	G_S = \langle g_i | g_i^2 = 1\rangle = \langle g_i | g_i^2 = 1\text{ and } g_i \text{ is non-trivial}\rangle
\]
where $i$ ranges from 1 to $d$. When $g_i$ is trivial the problem is essentially lower-dimensional
\begin{itemize}
\item[1)] If $S_i^{-1} = 0$ and $S_i^{1} = 0$, $S = S_i^0$ so the $x_i$ variable plays no role since $E_i S Q= SQ$.
\item[2)] If $S_i^{-1} = 0$ and $S_i^1 \neq 0$ then
\[
	E_i SQ = SQ - \ol{x}_i(S_i^{1}Q_i^{-1}+S_i^0 Q_i^{-1}) = SQ
\]
which is the same case as before
\item[3)] If $S_i^{-1}\neq 0$ but $S_i^{1} = 0$ then $E_iSQ$ is non-trivial but since we cannot take any steps away from the $x_i = 0$ hyperplane, we see that $S$ is essentially of lower dimension.
\end{itemize}
This allows us to define the \emph{dimension} of $G_S$ or of the random walk model itself. Let $\dim(S)$ be the number of non-trivial generators $g_i$ in this representation of the group $G_S$. Without loss of generality we will speak of models set in $d$ dimensions that have a full set of generators. Firstly, we can easily show that there are infinitely many different group sizes that are manifested among the different dimensions. Let $S$ be a $(d-1)$-dimensional model with group $G$ of largest finite order (number of elements) for that dimension. Set this model in $d$ dimensions and define a new model $S' = S(x_d+\ol{x}_d)$. This model has a non-trivial generator $g_d$ that $S$ lacks. Therefore the group $G'$ of this new model satisfies
\[
	G' \cong G \times \langle g_d \rangle
\]
and $|G'| = 2|G|$ and continuing in this manner allows us to write an infinite sequence of distinct groups. That is, not only are there infinitely many distinct groups but also, there is a group in $d$ dimensions not encountered before in $(d-1)$ dimensions (neither in order nor structure). For example the orders of the finite groups set in 2 dimensions are 1,2,4,6 and 8. So in 3 dimensions we have groups of order 1,2,4,6,8,12 and 16 at least. In fact Steve's paper shows that there are even more - of orders 24 and 48 too.

%========================================
\section{Coxeter Groups}
%========================================

This discussion would be short without delving into the subject of Coxeter groups. The $d$-dimensional groups defined above are Coxeter groups - a generalization of the dihedral groups encountered in dimension 2 random walk models. Coxeter groups are defined as
\[
	G = \langle r_1,\ldots, r_d |(r_ir_j)^{m_{i,j}} = 1\rangle
\]
where $m_{i,i} = 1$ and $m_{i,j} \ge 2$ for all $i\neq j$, and $m_{i,j}$ could also be infinite. We can call these $r$'s reflections. When $m_{i,j} = 2$ the reflections commute. Also $m_{i,j} = m_{j,i}$. The \emph{rank} of $G$ is the number of generators of in this specific representation of $G$. Often $G$ will have more than one representation with different numbers of generators. So in our discussion of groups for random walks with small steps the rank of $G$ is the dimension of the model. We can draw a graph for each Coxeter group
\begin{itemize}
\item label vertices with the generators of $G$
\item connect the vertices $g_i$ and $g_j$ if $m_{i,j} \ge 3$
\item label the edge with $m_{i,j}$ if $m_{i,j}\ge 4$.
\end{itemize}
If the graph has different components, the Coxeter group factors as a direct product of Coxeter groups represented by the individual components. We can see this in terms of the step polynomial $S$. Suppose that $S$ factors as
\[
	S = T(x_1,\ldots,x_k)R(x_{k+1},\ldots,x_d)
\]
For a fixed $1 \leq i \leq k$ we have
\[
	S_i^j = [x_i^j]S = R[x_i^j]T = R T_i^j
\]
and therefore
\[
	g_i(x_i) = \ol{x}_i\frac{T_i^{-1}R}{T_i^1R} = \ol{x}_i\frac{T_i^{-1}}{T_i^1}
\]
Similarly if $k+1\leq j\leq d$ then
\[
	g_j(x_j) = \ol{x}_j\frac{R_i^{-1}}{R_i^1}
\]
Clearly
\[
	g_ig_j = g_jg_i
\]
since they don't interfere with each other. Since $i$ and $j$ were arbitrary, this means that the generators from different sets commute and the Coxeter graph of the group factors into two pieces. We can write
\[
	G_S = G_T \times G_R
\]
where the subscripts highlight which parts of $S$ they groups are generated by. This factoring also allows us to write the generating function $Q$ as a Hadamard product of lower-dimensional generating functions. We can call each component of the graph a prime model. We can say that a simple small component is one whose group doesn't split non-trivially (its graph has one component).

It's worth mentioning that it doesn't seem that all Coxeter groups are manifested in random walk models. Looking at the classification of finite Coxeter groups, those of rank 2 have orders 4, 6, 8, 12, 10 and 2$p$ (manifested in a $p$-gon). Namely the groups of orders 10,12 and 2$p$ (for most $p$) are not manifested in random walk models with small steps on the quarter plane. It appears that the $p$-gon symmetries are not present for most $p$ because of the crystallographic nature of the problem. It's quite possible that the Crystallographic restriction theorem will apply to the random walk model groups in all dimensions. However, as can be seen, a lot of random walk model groups in 2 dimensions have generators that don't map the lattice back onto itself in a simple way because the transformations are bi-rational. We can define classes of groups.

\begin{itemize}
\item a model is "nice" iff $g_i(x_i) \in \ze[x_1,\ol{x}_1\ldots,x_d,\ol{x}_d]$
\item a model is "crystallographic" iff $g_i(x_1^{n_1}\ldots x_d^{n_d}) = x_1^{m_1}\ldots x_d^{m_d}$ for some integers $n_j$ and $m_j$.
\end{itemize}
The names are suggestive because the generators in crystallographic random walk models map monomials to monomials and we can see if $G$ is crystallographic. The next step up are "nice" models where monomials are mapped to Laurent polynomials. The rest of the classes have a generator mapping some monomials to rational functions. Perhaps those first two types of models can be represented via reflections and will thus have Weyl group.

%========================================
\section{Orbit Sums and Kernel Method}
%========================================

Since $G = G_S$ preserves the kernel $K$ we can construct an orbit sum where we have
\[
	\sum_{g\in G} \sign(g) g(x_1\ldots x_d Q) = \frac{1}{K(x_1,\ldots,x_d)}\sum_{g\in G} \sign(g) g(x_1\ldots x_d)
\]
We can write this equation more compactly via
\[
	D_G = \sum_{g \in G}\sign(g), \quad D_G' = D_g x_1\ldots x_d
\]
Then the identity reads in one of two followin ways
\[
	D_G'Q = \frac{1}{K}D_G'(1)
\]
or
\[
	D_G'(KQ - 1) = 0
\]
We'll prove the identity shortly. Let $g_k$ be one of the non-trivial generators. Then we can form the factor group $G/\langle g_k \rangle$ which naturally embeds into $G$ as $H$. Since $g_k$ is of order 2 $H$ is normal in $G$ and has two cosets - $H$ and $Hg_k$. To make the discussion neater let $e_g$ be the sign of $g$. This allows us to write $D_G$ in another way
\begin{align*}
	D_G =\sum_{g \in H} e_g g + \sum_{g \in Hg_k} e_g g\\
	&=\sum_{g \in H} e_g g + \sum_{g \in H} (-1)e_g g g_k\\
	&= \left(\sum_{g \in H} e_g g\right)(1-g_k)
	&= D_H(1-g_k)
\end{align*}
We can prove the identity by starting with the unlabeled functional equation written as
\[
	KQ + \sum_{\emptyset \subsetneq I \subseteq [d]}(-1)^{|I|}\ol{x}_{I} S_I^{-\one} Q_I^{\oo} = 1
\]
We see that
\[
	D_G'(\ol{x}_{I} S_I^{-\one} Q_I^{\oo}) = D_G(x_{[d]\backslash I} S_I^{-\one} Q_I^{\oo})
\]
Now, if we let $i \in I$ then $D_G = D_{H_i}(1-g_i)$ and since there are no variables $x_i$ in the expression $D_G$ klls that term. Applying $D_G'$ over the functional equation yields the result.

For a dihedral group of order 4 we have
\[
	D_G = 1- g_1 +g_2 g_1 -g_1g_2g_1 = (1+g_2)(1-g_1)
\]
but another way to write it is by using the relation
\[
	g_1g_2g_1g_2 = 1
\]
and multiplying both sides by $g_2$. We find that
\[
	D_G = 1 - g_1 -g_2 + g_2 g_1 = (1-g_2)(1-g_1)
\]

For a dihedral group of order 6 we have
\begin{align*}
	D_G &= 1- g_1 + g_2 g_1 -g_1 g_2 g_1+g_2 g_1 g_2 g_1 -g_1 g_2 g_1 g_2 g_1 \\
	&= (1+g_1 g_2+(g_1 g_2)^2)(1-g_1)\\
	&=(1+g_1g_2+g_2g_1)(1-g_1)
\end{align*}
The dihedral group of order 8 has 
\[
	D_G = (1+g_1g_2+g_2g_1+g_1g_2g_1g_2)(1-g_1)
\]
This also works for the group generated by two elements within a larger Coxeter group. We have 
\[
	D_{\langle g_i,g_j\rangle} = (1+g_ig_j+\cdots + (g_ig_j)^{m_{i,j}-1})(1-g_i)
\]
where we can exchange the meaning of $i$ and $j$ to get the alternative form. These calculations show that we probably won't be able to factor more than one factor $(1-g_k)$ out of the operator.
%========================================
\section{Half-Orbit Sums}
%========================================

Sometimes applying the operator $D_G'$ to 1 gives 0. In the cases studied it seems to happen when we have $(x,y) \mapsto (y,x)$ as part of the group. Let $h = h_{i,j}$ be the map $(x_i,x_j) \mapsto (x_j,x_i)$. We can form the quotient group $G/\langle h\rangle$ as before and find $H$ a subgroup of $G$ where $H \cong G/\langle h\rangle$. Since $h$ has order 2 we find that $H$ is normal and $G = H \cup Hh$. If $\sign(h) = -1$ we have
\[
	D_G = D_H (1-h),\quad (1-h)(x_1\ldots x_d) = 0
\]
and therefore this result is $D_G'(1) = 0$. This is an example where we would have to use the half-orbit method. But does this mean that every situation with $D_G'(1) = 0$ is a result of having an element like $h$ within $G$. One can certainly imagine higher order permutations being present within $G$ that can also lead to orbit of zero. Similarly, there may be group elements that are not permutations that also lead to orbit of zero. We can treat some of the possibilities as follows. Let $g\in G$ fix $x = x_1\ldots x_d$. That is 

%========================================
\section{Infinite Group Orbits}
%========================================

For an infinite group of dimension 2 the orbit operator will perhaps have the following form
\[
	D_G = (1+g+g^2+\cdots )(1-g_1) = (1-g)^{-1}(1-g_1)
\]
where $ g = g_1 g_2$ and we seek an inverse of $1-g$. However this will most likely be a difficult matter because $g$ can be rather complicated in some cases.
\endinput
